\documentclass[11pt]{article}
\usepackage[margin=18mm,headheight=22pt]{geometry}
\usepackage{amsmath}
\usepackage{fontspec}
\usepackage{unicode-math}
\setmainfont{NewCM10}[Extension=.otf,UprightFont=*-Book,BoldFont=*-Bold,ItalicFont=*-BookItalic,BoldItalicFont=*-BoldItalic]
\setsansfont{NewCMSans10}[Extension=.otf,UprightFont=*-Book,BoldFont=*-Bold,ItalicFont=*-BookOblique,BoldItalicFont=*-BoldOblique]
\setmathfont{NewCMMath-Book.otf}
\usepackage{xcolor}
\definecolor{FAUblue}{HTML}{04316A}
\definecolor{Paperblue}{HTML}{EEF2F7}
\definecolor{EEIgold}{HTML}{CC9C42}
\usepackage{etoolbox,enumitem,booktabs,tabularx,array,needspace,microtype}
\usepackage[most]{tcolorbox}
\usepackage{titlesec,fancyhdr}
\usepackage[hidelinks]{hyperref}
\usepackage{draftwatermark}
\SetWatermarkText{Made by Thanate Amornjiraporn (Focus)}
\SetWatermarkColor{black!7}\SetWatermarkFontSize{30pt}\SetWatermarkAngle{45}
\setlength{\parindent}{0pt}
\setlength{\parskip}{5pt}
\setlist{itemsep=2pt,topsep=3pt,leftmargin=1.6em,beginpenalty=10000,endpenalty=10000}
\raggedbottom
\clubpenalty=10000 \widowpenalty=10000 \displaywidowpenalty=10000 \interlinepenalty=10000
\setcounter{secnumdepth}{0}
\setcounter{tocdepth}{1}
\pagestyle{fancy}\fancyhf{}
\fancyhead[L]{\small\sffamily\color{FAUblue} From Fundamentals to AI Apocalypse · APDR · Exercise 1}
\fancyhead[R]{\small\sffamily Computer Exercise 1: Worked solutions}
\fancyfoot[C]{\small\thepage}
\titleformat{\section}{\Large\bfseries\sffamily\color{FAUblue}}{}{0pt}{}
\titleformat{\subsection}{\large\bfseries\sffamily\color{FAUblue}}{}{0pt}{}
\titlespacing*{\section}{0pt}{16pt}{6pt}
\preto\section{\Needspace{8\baselineskip}}
\preto\subsection{\Needspace{6\baselineskip}}
\newtcolorbox{equationcard}[1]{colback=Paperblue,colframe=Paperblue,boxrule=0pt,arc=0pt,left=10pt,right=10pt,top=6pt,bottom=6pt,title={#1},fonttitle=\small\bfseries\sffamily,coltitle=FAUblue,before skip=8pt,after skip=8pt}
\newtcolorbox{workedbox}[1]{colback=white,colframe=EEIgold,boxrule=0pt,leftrule=3pt,arc=0pt,left=10pt,top=4pt,bottom=4pt,title={#1},fonttitle=\small\bfseries\sffamily,coltitle=FAUblue,attach title to upper={\par},before skip=8pt,after skip=8pt}
\newtcolorbox{recallbox}[1]{colback=white,colframe=FAUblue,boxrule=0pt,leftrule=3pt,arc=0pt,left=10pt,top=4pt,bottom=4pt,title={#1},fonttitle=\small\bfseries\sffamily,coltitle=FAUblue,attach title to upper={\par},before skip=8pt,after skip=8pt}
\usepackage{environ}
\newsavebox\measurebox
% Typeset the section once invisibly to measure it; start a new page if it will not fit in what is left.
\makeatletter
\newif\ifkeep@measuring
\NewEnviron{keepsection}{%
  \ifkeep@measuring \BODY \else
    \edef\keep@sec{\the\value{section}}\edef\keep@sub{\the\value{subsection}}%
    \setbox\measurebox\vbox{\hsize\linewidth\keep@measuringtrue\def\Needspace##1{}\BODY}%
    \setcounter{section}{\keep@sec}\setcounter{subsection}{\keep@sub}%
    \dimen2=\ht\measurebox \advance\dimen2\dp\measurebox
    \dimen@=\textheight
    \ifdim\pagegoal<\maxdimen
      \dimen@=\pagegoal \advance\dimen@-\pagetotal
      \ifdim\dimen2>\dimen@ \clearpage \dimen@=\textheight \fi
    \fi
    % if it fits, nothing inside may force an early break
    \ifdim\dimen2>\dimen@ \BODY \else \begingroup\def\Needspace##1{}\BODY\endgroup \fi
  \fi}
\makeatother
\newtcolorbox{notebox}{colback=Paperblue,colframe=Paperblue,boxrule=0pt,arc=0pt,left=8pt,right=8pt,top=4pt,bottom=4pt,before skip=6pt,after skip=6pt}
\newtcolorbox{plainbox}{colback=white,colframe=FAUblue,boxrule=0pt,leftrule=3pt,arc=0pt,left=10pt,right=8pt,top=4pt,bottom=4pt,before skip=6pt,after skip=8pt}
\usetikzlibrary{arrows.meta}
\title{\sffamily\bfseries\color{FAUblue}Computer Exercise 1: Worked solutions}
\author{Algorithms, Programming, and Data Representation \ensuremath{\cdot} Module 92313 \ensuremath{\cdot} Exercise 1\\[2pt]\small From Fundamentals to AI Apocalypse · personal study guide}
\date{}
\begin{document}
\maketitle
\begingroup\footnotesize\setlength{\parskip}{0pt}\tableofcontents\endgroup
\clearpage

\Needspace{9\baselineskip}
{\large File signature forensics works from bytes and validated structures, never from filenames alone. Try each task before opening its solution. The ZIP family and the personalized sample offsets need the extra care explained below.}


\begin{keepsection}
\Needspace{5\baselineskip}
\section{Task 1: Read only the requested prefix}

\textbf{Connected topic:} Computers, programs, and data representation (Chapter 2)

\begin{keepsection}
\Needspace{7\baselineskip}
\subsection{1 \ensuremath{\cdot} read\_magic\_bytes(filepath, n=16)}

\Needspace{4\baselineskip}
Try first: decide what the function must return for a short file, an empty file, and n=0 before writing it.

\nopagebreak[4]{\small\itshape Answers: see “Task 1: Read only the requested prefix” under Answers at the end.}
\end{keepsection}
\end{keepsection}

\begin{keepsection}
\Needspace{5\baselineskip}
\section{Task 2: Create a useful signature table}

\textbf{Connected topic:} Integer representations and binary data (§3.2)

\begin{keepsection}
\Needspace{7\baselineskip}
\subsection{2 \ensuremath{\cdot} SIGNATURES}

\Needspace{4\baselineskip}
Try first: list real leading byte patterns for distinct formats, then separate container formats from their packaged subtypes.

\nopagebreak[4]{\small\itshape Answers: see “Task 2: Create a useful signature table” under Answers at the end.}
\end{keepsection}
\end{keepsection}

\begin{keepsection}
\Needspace{5\baselineskip}
\section{Task 3: Identify a type and resolve ZIP subtypes}

\textbf{Connected topic:} Integer representations and binary data (§3.2)

\begin{keepsection}
\Needspace{6\baselineskip}
\subsection{3 \ensuremath{\cdot} identify\_file\_type(magic)}

\Needspace{3\baselineskip}
Try first: separate “what the first bytes suggest” from “what the whole container contains.”

\nopagebreak[4]{\small\itshape Answers: see “Task 3: Identify a type and resolve ZIP subtypes” under Answers at the end.}
\end{keepsection}
\end{keepsection}

\begin{keepsection}
\Needspace{5\baselineskip}
\section{Task 4: Hash a file and verify integrity}

\textbf{Connected topic:} Computers, programs, and data representation (Chapter 2)

\begin{keepsection}
\Needspace{7\baselineskip}
\subsection{4 \ensuremath{\cdot} compute\_checksum / verify\_integrity}

\Needspace{4\baselineskip}
Try first: process the file in fixed-size chunks so memory use does not grow with the file size.

\nopagebreak[4]{\small\itshape Answers: see “Task 4: Hash a file and verify integrity” under Answers at the end.}
\end{keepsection}
\end{keepsection}

\begin{keepsection}
\Needspace{5\baselineskip}
\section{Task 5: Find mismatched file extensions recursively}

\textbf{Connected topic:} Computers, programs, and data representation (Chapter 2)

\begin{keepsection}
\Needspace{7\baselineskip}
\subsection{5 \ensuremath{\cdot} find\_mismatched\_extensions(root\_dir)}

\Needspace{4\baselineskip}
Try first: compare the suffix claim with the type discovered from the file bytes and, for ZIPs, its member list.

\nopagebreak[4]{\small\itshape Answers: see “Task 5: Find mismatched file extensions recursively” under Answers at the end.}
\end{keepsection}
\end{keepsection}

\begin{keepsection}
\Needspace{5\baselineskip}
\section{Task 6: Find files that validly match multiple formats}

\textbf{Connected topic:} Integer representations and binary data (§3.2)

\begin{keepsection}
\Needspace{7\baselineskip}
\subsection{6 \ensuremath{\cdot} detect\_polyglot(filepath)}

\Needspace{4\baselineskip}
Try first: collect and validate candidates independently instead of returning the first signature hit.

\nopagebreak[4]{\small\itshape Answers: see “Task 6: Find files that validly match multiple formats” under Answers at the end.}
\end{keepsection}
\end{keepsection}

\begin{keepsection}
\Needspace{5\baselineskip}
\section{Task 7: Format an exact hexdump}

\textbf{Connected topic:} Character and byte representations (§3.3)

\begin{keepsection}
\Needspace{6\baselineskip}
\subsection{7 \ensuremath{\cdot} bytes\_to\_hexdump(data, width=16)}

\Needspace{3\baselineskip}
Try first: work from chunks so offsets and columns are stable even on the last, partial line.

\nopagebreak[4]{\small\itshape Answers: see “Task 7: Format an exact hexdump” under Answers at the end.}
\end{keepsection}
\end{keepsection}

\begin{keepsection}
\Needspace{5\baselineskip}
\section{Task 8: Classify common text encodings from bytes}

\textbf{Connected topic:} Character and byte representations (§3.3)

\begin{keepsection}
\Needspace{6\baselineskip}
\subsection{8 \ensuremath{\cdot} detect\_text\_encoding(filepath)}

\Needspace{3\baselineskip}
Try first: check BOMs first, then use byte patterns and a strict UTF-8 validity check.

\nopagebreak[4]{\small\itshape Answers: see “Task 8: Classify common text encodings from bytes” under Answers at the end.}
\end{keepsection}
\end{keepsection}

\begin{keepsection}
\Needspace{6\baselineskip}
\section{Task 9: Find the first byte after the valid file structure}

\Needspace{3\baselineskip}
\textbf{Connected topic:} Integer representations and binary data (§3.2)

\begin{keepsection}
\Needspace{7\baselineskip}
\subsection{9 \ensuremath{\cdot} find\_appended\_data(filepath, filetype)}

\Needspace{4\baselineskip}
Try first: parse the end of the format structure and compare that offset with the full file length.

\nopagebreak[4]{\small\itshape Answers: see “Task 9: Find the first byte after the valid file structure” under Answers at the end.}
\end{keepsection}
\end{keepsection}

\begin{keepsection}
\Needspace{5\baselineskip}
\section{Task 10: Aggregate the forensic report and write the personalized justification}

\textbf{Connected topic:} Computers, programs, and data representation (Chapter 2)

\begin{keepsection}
\Needspace{7\baselineskip}
\subsection{10 \ensuremath{\cdot} generate\_forensics\_report(root\_dir)}

\Needspace{4\baselineskip}
Try first: reuse one read per file, make a JSON-safe record, then measure the runtime on 500 or more files.

\nopagebreak[4]{\small\itshape Answers: see “Task 10: Aggregate the forensic report and write the personalized justification” under Answers at the end.}
\end{keepsection}

\begin{keepsection}
\Needspace{7\baselineskip}
\subsection{Written justification \ensuremath{\cdot} use your own sample offsets}

\Needspace{4\baselineskip}
Try first: record the polyglot file and appended-data boundary from your assigned samples; the seeded filenames and offsets differ by learner.

\nopagebreak[4]{\small\itshape Answers: see “Task 10: Aggregate the forensic report and write the personalized justification” under Answers at the end.}
\end{keepsection}
\end{keepsection}


\clearpage
\section{Answers}
\begin{keepsection}\subsection{Task 1: Read only the requested prefix}

\begin{workedbox}{1 \ensuremath{\cdot} read\_magic\_bytes(filepath, n=16) \ensuremath{\cdot} worked solution}
\Needspace{11\baselineskip}
\begin{enumerate}
\item The input is a path to an arbitrary binary file. Open it in binary mode so the bytes are not decoded or translated. The result is the first n bytes, not the whole file.
\item A bounded read naturally returns fewer than n bytes at end-of-file. Reject a negative n: Python treats \texttt{read(-1)} as “read the rest”, which would break the memory limit.
\item Close the handle on success and on exceptions with a context manager. Check a short file and n=0: the returned values should be the bytes available and \texttt{b""}, respectively.
\item \Needspace{3\baselineskip}
\textbf{The whole function:}
\begin{tcolorbox}[colback=Paperblue,colframe=Paperblue,boxrule=0pt,arc=0pt,left=7pt,right=7pt,top=5pt,bottom=5pt]\small\sloppy
\noindent\hspace*{0em}\texttt{def read\_magic\_bytes(filepath, n=16):}\par
\noindent\hspace*{1em}\texttt{if n \ensuremath{<} 0:}\par
\noindent\hspace*{3em}\texttt{raise ValueError("n must be nonnegative")}\par
\noindent\hspace*{1em}\texttt{with open(filepath, "rb") as f:}\par
\noindent\hspace*{3em}\texttt{return f.read(n)}\par
\end{tcolorbox}
\end{enumerate}
\end{workedbox}

\end{keepsection}
\begin{keepsection}\subsection{Task 2: Create a useful signature table}

\begin{workedbox}{2 \ensuremath{\cdot} SIGNATURES \ensuremath{\cdot} worked solution}
\Needspace{16\baselineskip}
\begin{enumerate}
\item A signature is bytes, so store bytes literals and match from offset zero. Include multiple common JPEG starts and both GIF header versions. Add at least twelve distinct formats.
\item DOCX, XLSX, PPTX, and JAR are ZIP containers: their local-file headers share the same PK prefix. Repeating that one prefix under four names does not distinguish them. Add these names as archive subtypes in a second-stage classifier, not as unique magic signatures.
\item The sample table below covers thirteen signatures. It is a candidate table; a matching prefix is evidence to validate the format structure, not a complete proof that the file is valid.
\item \Needspace{3\baselineskip}
\textbf{The table:}
\begin{tcolorbox}[colback=Paperblue,colframe=Paperblue,boxrule=0pt,arc=0pt,left=7pt,right=7pt,top=5pt,bottom=5pt]\small\sloppy
\noindent\hspace*{0em}\texttt{SIGNATURES = \{}\par
\noindent\hspace*{1em}\texttt{"PNG": (b"\textbackslash{}x89PNG\textbackslash{}r\textbackslash{}n\textbackslash{}x1a\textbackslash{}n",),}\par
\noindent\hspace*{1em}\texttt{"JPEG": (b"\textbackslash{}xff\textbackslash{}xd8\textbackslash{}xff\textbackslash{}xe0", b"\textbackslash{}xff\textbackslash{}xd8\textbackslash{}xff\textbackslash{}xe1", b"\textbackslash{}xff\textbackslash{}xd8\textbackslash{}xff\textbackslash{}xdb"),}\par
\noindent\hspace*{1em}\texttt{"GIF": (b"GIF87a", b"GIF89a"),}\par
\noindent\hspace*{1em}\texttt{"PDF": (b"\%PDF-",),}\par
\noindent\hspace*{1em}\texttt{"ZIP": (b"PK\textbackslash{}x03\textbackslash{}x04", b"PK\textbackslash{}x05\textbackslash{}x06", b"PK\textbackslash{}x07\textbackslash{}x08"),}\par
\noindent\hspace*{1em}\texttt{"GZIP": (b"\textbackslash{}x1f\textbackslash{}x8b",),}\par
\noindent\hspace*{1em}\texttt{"ELF": (b"\textbackslash{}x7fELF",),}\par
\noindent\hspace*{1em}\texttt{"RTF": (b"\{\textbackslash{}\textbackslash{}rtf",),}\par
\noindent\hspace*{1em}\texttt{"BMP": (b"BM",),}\par
\noindent\hspace*{1em}\texttt{"SQLite": (b"SQLite format 3\textbackslash{}x00",),}\par
\noindent\hspace*{1em}\texttt{"7z": (b"7z\textbackslash{}xbc\textbackslash{}xaf\textbackslash{}x27\textbackslash{}x1c",),}\par
\noindent\hspace*{1em}\texttt{"Ogg": (b"OggS",),}\par
\noindent\hspace*{1em}\texttt{"MP3": (b"ID3",),}\par
\noindent\hspace*{0em}\texttt{\}}\par
\end{tcolorbox}
\end{enumerate}
\end{workedbox}

\end{keepsection}
\begin{keepsection}\subsection{Task 3: Identify a type and resolve ZIP subtypes}

\begin{workedbox}{3 \ensuremath{\cdot} identify\_file\_type(magic) \ensuremath{\cdot} worked solution}
\Needspace{22\baselineskip}
\begin{enumerate}
\item For ordinary types, compare the available prefix to each signature; \texttt{startswith} is safe for short byte strings and cannot index past the end. If several patterns match, prefer the longest matching signature.
\item The sheet’s request to distinguish DOCX/XLSX/PPTX/JAR from magic alone is underspecified: their first bytes are shared. With only the first 16 bytes, the correct answer is ZIP. To refine it, inspect the archive directory: \texttt{word/document.xml} indicates DOCX, \texttt{xl/workbook.xml} XLSX, \texttt{ppt/presentation.xml} PPTX, and \texttt{META-INF/MANIFEST.MF} JAR.
\item Use “unknown” only when no signature matches. Keep the container decision explicit so a generic ZIP is not mislabeled as an office document.
\item \Needspace{3\baselineskip}
\textbf{The functions} (with \texttt{import zipfile} at the top of the file):
\begin{tcolorbox}[colback=Paperblue,colframe=Paperblue,boxrule=0pt,arc=0pt,left=7pt,right=7pt,top=5pt,bottom=5pt]\small\sloppy
\noindent\hspace*{0em}\texttt{def identify\_file\_type(magic):}\par
\noindent\hspace*{1em}\texttt{matches = [(len(sig), kind)}\par
\noindent\hspace*{6em}\texttt{for kind, signatures in SIGNATURES.items()}\par
\noindent\hspace*{6em}\texttt{for sig in signatures if magic.startswith(sig)]}\par
\noindent\hspace*{1em}\texttt{return max(matches)[1] if matches else "unknown"}\par
\noindent\hspace*{0em}\texttt{}\par
\noindent\hspace*{0em}\texttt{def identify\_zip\_members(source):}\par
\noindent\hspace*{1em}\texttt{with zipfile.ZipFile(source) as archive:}\par
\noindent\hspace*{3em}\texttt{names = set(archive.namelist())}\par
\noindent\hspace*{1em}\texttt{if "word/document.xml" in names:}\par
\noindent\hspace*{3em}\texttt{return "DOCX"}\par
\noindent\hspace*{1em}\texttt{if "xl/workbook.xml" in names:}\par
\noindent\hspace*{3em}\texttt{return "XLSX"}\par
\noindent\hspace*{1em}\texttt{if "ppt/presentation.xml" in names:}\par
\noindent\hspace*{3em}\texttt{return "PPTX"}\par
\noindent\hspace*{1em}\texttt{if "META-INF/MANIFEST.MF" in names:}\par
\noindent\hspace*{3em}\texttt{return "JAR"}\par
\noindent\hspace*{1em}\texttt{return "ZIP"}\par
\noindent\hspace*{0em}\texttt{}\par
\noindent\hspace*{0em}\texttt{def classify\_path(path):}\par
\noindent\hspace*{1em}\texttt{magic = read\_magic\_bytes(path, n=4096)}\par
\noindent\hspace*{1em}\texttt{kind = identify\_file\_type(magic)}\par
\noindent\hspace*{1em}\texttt{if kind == "ZIP":}\par
\noindent\hspace*{3em}\texttt{return identify\_zip\_members(path) if zipfile.is\_zipfile(path) else "unknown"}\par
\noindent\hspace*{1em}\texttt{return kind}\par
\end{tcolorbox}
\end{enumerate}
\end{workedbox}

\end{keepsection}
\begin{keepsection}\subsection{Task 4: Hash a file and verify integrity}

\begin{workedbox}{4 \ensuremath{\cdot} compute\_checksum / verify\_integrity \ensuremath{\cdot} worked solution}
\Needspace{15\baselineskip}
\begin{enumerate}
\item Create the requested hashlib digest by algorithm name. Unsupported names should raise the documented \texttt{ValueError} rather than silently switching algorithms.
\item Read a chunk, update the digest, and repeat until an empty bytes value signals EOF. \texttt{hexdigest} gives the stable printable representation.
\item Verification means “recompute and compare with the expected digest”; a match detects accidental or malicious changes only when the expected digest is trusted. It does not prove who created the file.
\item \Needspace{3\baselineskip}
\textbf{Both functions} (with \texttt{import hashlib, hmac}):
\begin{tcolorbox}[colback=Paperblue,colframe=Paperblue,boxrule=0pt,arc=0pt,left=7pt,right=7pt,top=5pt,bottom=5pt]\small\sloppy
\noindent\hspace*{0em}\texttt{def compute\_checksum(filepath, algo="sha256"):}\par
\noindent\hspace*{1em}\texttt{digest = hashlib.new(algo)}\par
\noindent\hspace*{1em}\texttt{with open(filepath, "rb") as f:}\par
\noindent\hspace*{3em}\texttt{for chunk in iter(lambda: f.read(65536), b""):}\par
\noindent\hspace*{5em}\texttt{digest.update(chunk)}\par
\noindent\hspace*{1em}\texttt{return digest.hexdigest()}\par
\noindent\hspace*{0em}\texttt{}\par
\noindent\hspace*{0em}\texttt{def verify\_integrity(filepath, expected, algo="sha256"):}\par
\noindent\hspace*{1em}\texttt{actual = compute\_checksum(filepath, algo)}\par
\noindent\hspace*{1em}\texttt{return hmac.compare\_digest(actual, expected.strip().lower())}\par
\end{tcolorbox}
\end{enumerate}
\end{workedbox}

\end{keepsection}
\begin{keepsection}\subsection{Task 5: Find mismatched file extensions recursively}

\begin{workedbox}{5 \ensuremath{\cdot} find\_mismatched\_extensions(root\_dir) \ensuremath{\cdot} worked solution}
\Needspace{21\baselineskip}
\begin{enumerate}
\item \texttt{os.walk} visits nested directories without a fixed depth. Normalize the suffix to lowercase and use a table; keep an unrecognized extension as "unknown" rather than guessing.
\item Classify content independently of the suffix. For a ZIP signature, open the archive and use Task 3’s member-name check to distinguish a plain ZIP from an office package or JAR.
\item Only after a signature lookup returns unknown, classify text from the bytes; this must work for a text file with any suffix. Then compare the independently detected actual type with the suffix claim. Test a nested wrong extension, uppercase suffix, missing suffix, and unknown actual signature.
\item \Needspace{4\baselineskip}
\textbf{The function.} It uses \texttt{EXTENSION\_TYPES}, a dictionary from lowercase suffix to type name such as \texttt{\{".png": "PNG", ".txt": "TXT"\}}, and the helpers from Tasks 3 and 8:
\begin{tcolorbox}[colback=Paperblue,colframe=Paperblue,boxrule=0pt,arc=0pt,left=7pt,right=7pt,top=5pt,bottom=5pt]\small\sloppy
\noindent\hspace*{0em}\texttt{def find\_mismatched\_extensions(root\_dir):}\par
\noindent\hspace*{1em}\texttt{mismatches = []}\par
\noindent\hspace*{1em}\texttt{for folder, \_, filenames in os.walk(root\_dir):}\par
\noindent\hspace*{3em}\texttt{for filename in filenames:}\par
\noindent\hspace*{5em}\texttt{path = os.path.join(folder, filename)}\par
\noindent\hspace*{5em}\texttt{extension = os.path.splitext(filename)[1].lower()}\par
\noindent\hspace*{5em}\texttt{claimed = EXTENSION\_TYPES.get(extension, "unknown")}\par
\noindent\hspace*{5em}\texttt{actual = classify\_path(path)}\par
\noindent\hspace*{5em}\texttt{if actual == "unknown" and detect\_text\_encoding(path) != "unknown":}\par
\noindent\hspace*{7em}\texttt{actual = "TXT" \# infer text from bytes, regardless of suffix}\par
\noindent\hspace*{5em}\texttt{if claimed != actual:}\par
\noindent\hspace*{7em}\texttt{mismatches.append(\{"path": path,}\par
\noindent\hspace*{15em}\texttt{"claimed\_type": claimed,}\par
\noindent\hspace*{15em}\texttt{"actual\_type": actual\})}\par
\noindent\hspace*{1em}\texttt{return mismatches}\par
\end{tcolorbox}
\end{enumerate}
\end{workedbox}

\end{keepsection}
\begin{keepsection}\subsection{Task 6: Find files that validly match multiple formats}

\begin{workedbox}{6 \ensuremath{\cdot} detect\_polyglot(filepath) \ensuremath{\cdot} worked solution}
\Needspace{17\baselineskip}
\begin{enumerate}
\item Load or stream the file once, then ask each supported format validator whether the bytes satisfy that format’s structure. Add each valid type once.
\item For a ZIP-based file, verify the central directory and inspect members. A string like \texttt{\%PDF-} somewhere inside a PNG is not, by itself, proof of a valid PDF; validate the PDF header position and its required ending structure too.
\item Return an empty list for no recognized valid format and a list of all valid matches otherwise. More than one match is the polyglot signal. Sort the result if deterministic output is required.
\item \Needspace{4\baselineskip}
\textbf{The function.} Each entry of \texttt{FORMAT\_VALIDATORS} maps a type name to a function that takes the file's bytes and returns whether they form a complete, valid file of that type:
\begin{tcolorbox}[colback=Paperblue,colframe=Paperblue,boxrule=0pt,arc=0pt,left=7pt,right=7pt,top=5pt,bottom=5pt]\small\sloppy
\noindent\hspace*{0em}\texttt{def detect\_polyglot(filepath):}\par
\noindent\hspace*{1em}\texttt{if isinstance(filepath, bytes):}\par
\noindent\hspace*{3em}\texttt{data = filepath}\par
\noindent\hspace*{1em}\texttt{else:}\par
\noindent\hspace*{3em}\texttt{with open(filepath, "rb") as f:}\par
\noindent\hspace*{5em}\texttt{data = f.read()}\par
\noindent\hspace*{1em}\texttt{matches = set()}\par
\noindent\hspace*{1em}\texttt{for kind, validator in FORMAT\_VALIDATORS.items():}\par
\noindent\hspace*{3em}\texttt{if validator(data): \# each validator checks the complete format structure}\par
\noindent\hspace*{5em}\texttt{matches.add(kind)}\par
\noindent\hspace*{1em}\texttt{return sorted(matches)}\par
\end{tcolorbox}
\end{enumerate}
\end{workedbox}

\end{keepsection}
\begin{keepsection}\subsection{Task 7: Format an exact hexdump}

\begin{workedbox}{7 \ensuremath{\cdot} bytes\_to\_hexdump(data, width=16) \ensuremath{\cdot} worked solution}
\Needspace{15\baselineskip}
\begin{enumerate}
\item For each offset 0, width, 2\ensuremath{\times}width, …, slice at most width bytes. Format the offset as eight uppercase hex digits and each byte as two uppercase hex digits.
\item Pad the hex column to width \ensuremath{\times} 3 characters. This makes the ASCII bar start in one column even when the last chunk is short.
\item Printable ASCII is byte values 0x20 through 0x7E inclusive; render every other value as a dot. Keep the separator, padding, and no-final-newline behavior exact.
\item \Needspace{3\baselineskip}
\textbf{The whole function:}
\begin{tcolorbox}[colback=Paperblue,colframe=Paperblue,boxrule=0pt,arc=0pt,left=7pt,right=7pt,top=5pt,bottom=5pt]\small\sloppy
\noindent\hspace*{0em}\texttt{def bytes\_to\_hexdump(data, width=16):}\par
\noindent\hspace*{1em}\texttt{if width \ensuremath{<}= 0:}\par
\noindent\hspace*{3em}\texttt{raise ValueError("width must be positive")}\par
\noindent\hspace*{1em}\texttt{lines = []}\par
\noindent\hspace*{1em}\texttt{for offset in range(0, len(data), width):}\par
\noindent\hspace*{3em}\texttt{chunk = data[offset:offset + width]}\par
\noindent\hspace*{3em}\texttt{hex\_column = " ".join(f"\{byte:02X\}" for byte in chunk)}\par
\noindent\hspace*{3em}\texttt{hex\_column = hex\_column.ljust(width * 3)}\par
\noindent\hspace*{3em}\texttt{ascii\_column = "".join(chr(b) if 0x20 \ensuremath{<}= b \ensuremath{<}= 0x7E else "."}\par
\noindent\hspace*{13em}\texttt{for b in chunk)}\par
\noindent\hspace*{3em}\texttt{lines.append(f"\{offset:08X\} \{hex\_column\} |\{ascii\_column\}|")}\par
\noindent\hspace*{1em}\texttt{return "\textbackslash{}n".join(lines)}\par
\end{tcolorbox}
\end{enumerate}
\end{workedbox}

\end{keepsection}
\begin{keepsection}\subsection{Task 8: Classify common text encodings from bytes}

\begin{workedbox}{8 \ensuremath{\cdot} detect\_text\_encoding(filepath) \ensuremath{\cdot} worked solution}
\Needspace{27\baselineskip}
\begin{enumerate}
\item A UTF-8 BOM is EF BB BF; UTF-16 little-endian is FF FE; UTF-16 big-endian is FE FF. BOM checks take priority over heuristics.
\item Without a BOM, UTF-16 text often has many zero bytes in one parity: ASCII-heavy LE text has zeros at odd indexes, while BE has zeros at even indexes. This heuristic needs enough bytes to show the pattern; short no-BOM strings are ambiguous and should be reported as unknown when they contain zero bytes. Binary data can also imitate the pattern.
\item ASCII is a subset of UTF-8, so test all-ASCII bytes first when the result labels distinguish them. For non-ASCII bytes, try strict UTF-8 decoding explicitly; if it fails, use the assignment’s Latin-1 fallback. This is explicit validation, not automatic encoding detection.
\item \Needspace{4\baselineskip}
\textbf{The function,} written on bytes so Task 10 can reuse it; \texttt{detect\_text\_encoding(filepath)} only reads the file and passes the bytes on:
\begin{tcolorbox}[colback=Paperblue,colframe=Paperblue,boxrule=0pt,arc=0pt,left=7pt,right=7pt,top=5pt,bottom=5pt]\small\sloppy
\noindent\hspace*{0em}\texttt{def detect\_text\_encoding(filepath):}\par
\noindent\hspace*{1em}\texttt{with open(filepath, "rb") as f:}\par
\noindent\hspace*{3em}\texttt{return detect\_text\_encoding\_bytes(f.read())}\par
\noindent\hspace*{0em}\texttt{}\par
\noindent\hspace*{0em}\texttt{def detect\_text\_encoding\_bytes(data):}\par
\noindent\hspace*{1em}\texttt{if data.startswith(b"\textbackslash{}xef\textbackslash{}xbb\textbackslash{}xbf"):}\par
\noindent\hspace*{3em}\texttt{return "UTF-8"}\par
\noindent\hspace*{1em}\texttt{if data.startswith(b"\textbackslash{}xff\textbackslash{}xfe"):}\par
\noindent\hspace*{3em}\texttt{return "UTF-16-LE"}\par
\noindent\hspace*{1em}\texttt{if data.startswith(b"\textbackslash{}xfe\textbackslash{}xff"):}\par
\noindent\hspace*{3em}\texttt{return "UTF-16-BE"}\par
\noindent\hspace*{1em}\texttt{paired = data[:len(data) - (len(data) \% 2)]}\par
\noindent\hspace*{1em}\texttt{even\_zeros = sum(paired[i] == 0 for i in range(0, len(paired), 2))}\par
\noindent\hspace*{1em}\texttt{odd\_zeros = sum(paired[i] == 0 for i in range(1, len(paired), 2))}\par
\noindent\hspace*{1em}\texttt{if len(paired) \ensuremath{>}= 8 and odd\_zeros \ensuremath{>} len(paired) // 4 and odd\_zeros \ensuremath{>} even\_zeros:}\par
\noindent\hspace*{3em}\texttt{return "UTF-16-LE"}\par
\noindent\hspace*{1em}\texttt{if len(paired) \ensuremath{>}= 8 and even\_zeros \ensuremath{>} len(paired) // 4 and even\_zeros \ensuremath{>} odd\_zeros:}\par
\noindent\hspace*{3em}\texttt{return "UTF-16-BE"}\par
\noindent\hspace*{1em}\texttt{if b"\textbackslash{}x00" in data:}\par
\noindent\hspace*{3em}\texttt{return "unknown" \# short no-BOM UTF-16 and binary data are ambiguous}\par
\noindent\hspace*{1em}\texttt{controls = sum(byte \ensuremath{<} 0x20 and byte not in (9, 10, 12, 13) for byte in data)}\par
\noindent\hspace*{1em}\texttt{if data and controls \ensuremath{>} len(data) // 100:}\par
\noindent\hspace*{3em}\texttt{return "unknown" \# reject mostly binary/control content}\par
\noindent\hspace*{1em}\texttt{if not data or all(byte \ensuremath{<} 0x80 for byte in data):}\par
\noindent\hspace*{3em}\texttt{return "ASCII"}\par
\noindent\hspace*{1em}\texttt{try:}\par
\noindent\hspace*{3em}\texttt{data.decode("utf-8", errors="strict")}\par
\noindent\hspace*{3em}\texttt{return "UTF-8"}\par
\noindent\hspace*{1em}\texttt{except UnicodeDecodeError:}\par
\noindent\hspace*{3em}\texttt{return "Latin-1"}\par
\end{tcolorbox}
\end{enumerate}
\end{workedbox}

\end{keepsection}
\begin{keepsection}\subsection{Task 9: Find the first byte after the valid file structure}

\begin{workedbox}{9 \ensuremath{\cdot} find\_appended\_data(filepath, filetype) \ensuremath{\cdot} worked solution}
\Needspace{27\baselineskip}
\begin{enumerate}
\item For PNG, start after its 8-byte signature. Each chunk is 4-byte big-endian length, 4-byte type, data, and 4-byte CRC; advance by 12 + length until a valid zero-length \texttt{IEND} chunk. The first byte after \texttt{IEND} is the candidate payload start.
\item For JPEG, find the EOI marker FF D9 after correctly handling marker segments and byte-stuffed scan data. For GIF, parse image/extension sub-blocks to the trailer byte. For PDF, use the final \texttt{\%\%EOF} marker and account for allowed trailing whitespace.
\item For ZIP, locate and validate the EOCD record; its 16-bit comment length determines the archive end. BMP stores its declared file size. RTF needs escaped-brace and binary-section handling; GZIP must distinguish concatenated members; ELF/SQLite derive the end from validated tables or page counts.
\item Return \texttt{None} when the structure is invalid, truncated, unsupported, or ends exactly at EOF. Marker search alone can match bytes inside payloads, so validate lengths and bounds before accepting an offset.
\item \Needspace{5\baselineskip}
\textbf{The PNG branch.} Write the parser on bytes as well: \texttt{find\_appended\_data(filepath, filetype)} reads the file once and calls \texttt{find\_appended\_data\_bytes(data, filetype)}, which runs the branch for that type (this is the PNG one):
\begin{tcolorbox}[colback=Paperblue,colframe=Paperblue,boxrule=0pt,arc=0pt,left=7pt,right=7pt,top=5pt,bottom=5pt]\small\sloppy
\noindent\hspace*{0em}\texttt{\# inside find\_appended\_data\_bytes(data, filetype), when filetype == "PNG"}\par
\noindent\hspace*{0em}\texttt{if not data.startswith(SIGNATURES["PNG"][0]):}\par
\noindent\hspace*{1em}\texttt{return None}\par
\noindent\hspace*{0em}\texttt{pos = 8 \# after the PNG signature}\par
\noindent\hspace*{0em}\texttt{while pos + 12 \ensuremath{<}= len(data):}\par
\noindent\hspace*{1em}\texttt{length = int.from\_bytes(data[pos:pos + 4], "big")}\par
\noindent\hspace*{1em}\texttt{kind = data[pos + 4:pos + 8]}\par
\noindent\hspace*{1em}\texttt{end = pos + 12 + length \# header + data + CRC}\par
\noindent\hspace*{1em}\texttt{if end \ensuremath{>} len(data):}\par
\noindent\hspace*{3em}\texttt{return None \# truncated chunk}\par
\noindent\hspace*{1em}\texttt{chunk = data[pos + 8:pos + 8 + length]}\par
\noindent\hspace*{1em}\texttt{stored\_crc = int.from\_bytes(data[pos + 8 + length:end], "big")}\par
\noindent\hspace*{1em}\texttt{if zlib.crc32(kind + chunk) \& 0xFFFFFFFF != stored\_crc:}\par
\noindent\hspace*{3em}\texttt{return None}\par
\noindent\hspace*{1em}\texttt{if kind == b"IEND":}\par
\noindent\hspace*{3em}\texttt{if length != 0:}\par
\noindent\hspace*{5em}\texttt{return None}\par
\noindent\hspace*{3em}\texttt{return None if end == len(data) else end}\par
\noindent\hspace*{1em}\texttt{pos = end}\par
\noindent\hspace*{0em}\texttt{return None}\par
\end{tcolorbox}
\end{enumerate}
\end{workedbox}

\end{keepsection}
\begin{keepsection}\subsection{Task 10: Aggregate the forensic report and write the personalized justification}

\begin{workedbox}{10 \ensuremath{\cdot} generate\_forensics\_report(root\_dir) \ensuremath{\cdot} worked solution}
\Needspace{31\baselineskip}
\begin{enumerate}
\item Walk every path once. Read that file’s bytes once, then pass the in-memory bytes to the signature, encoding, hexdump, polyglot, and appended-data helpers. Hash the same bytes with \texttt{hashlib.sha256}; do not reopen the path for each task.
\item For every file, record its path, first bytes, detected type, extension claim/mismatch, checksum, encoding result, polyglot types, appended offsets, and any bounded hexdump preview. If no known signature matches but the bytes pass Task 8 text checks, call the actual type \texttt{TXT} independent of its suffix. Keep values to JSON-compatible strings, numbers, lists, dictionaries, booleans, and nulls.
\item Keep aggregate mismatches in their own list if that is the required schema. Benchmark with \texttt{time.perf\_counter} on at least 500 files; report the measured elapsed time and environment rather than inventing a “within budget” result.
\item \Needspace{4\baselineskip}
\textbf{The loop} (with \texttt{import hashlib, io, os, zipfile}), the body of \texttt{generate\_forensics\_report(root\_dir)}:
\begin{tcolorbox}[colback=Paperblue,colframe=Paperblue,boxrule=0pt,arc=0pt,left=7pt,right=7pt,top=5pt,bottom=5pt]\small\sloppy
\noindent\hspace*{0em}\texttt{report = \{"mismatches": [], "files": []\}}\par
\noindent\hspace*{0em}\texttt{for folder, directories, names in os.walk(root\_dir):}\par
\noindent\hspace*{1em}\texttt{directories.sort()}\par
\noindent\hspace*{1em}\texttt{for name in sorted(names):}\par
\noindent\hspace*{3em}\texttt{path = os.path.join(folder, name)}\par
\noindent\hspace*{3em}\texttt{with open(path, "rb") as f:}\par
\noindent\hspace*{5em}\texttt{data = f.read() \# reuse these bytes; do not reopen for each task}\par
\noindent\hspace*{3em}\texttt{claimed = EXTENSION\_TYPES.get(os.path.splitext(name)[1].lower(), "unknown")}\par
\noindent\hspace*{3em}\texttt{actual = identify\_file\_type(data[:16])}\par
\noindent\hspace*{3em}\texttt{if actual == "ZIP" and zipfile.is\_zipfile(io.BytesIO(data)):}\par
\noindent\hspace*{5em}\texttt{actual = identify\_zip\_members(io.BytesIO(data))}\par
\noindent\hspace*{3em}\texttt{if actual == "unknown" and detect\_text\_encoding\_bytes(data) != "unknown":}\par
\noindent\hspace*{5em}\texttt{actual = "TXT" \# infer from bytes, never from the suffix}\par
\noindent\hspace*{3em}\texttt{types = detect\_polyglot(data)}\par
\noindent\hspace*{3em}\texttt{appended = \{kind: offset for kind in types}\par
\noindent\hspace*{9em}\texttt{if (offset := find\_appended\_data\_bytes(data, kind)) is not None\}}\par
\noindent\hspace*{3em}\texttt{row = \{"path": path, "magic\_hex": data[:16].hex(),}\par
\noindent\hspace*{6em}\texttt{"actual\_type": actual, "claimed\_type": claimed,}\par
\noindent\hspace*{6em}\texttt{"sha256": hashlib.sha256(data).hexdigest(),}\par
\noindent\hspace*{6em}\texttt{"text\_encoding": detect\_text\_encoding\_bytes(data),}\par
\noindent\hspace*{6em}\texttt{"polyglot\_types": types, "appended\_offsets": appended,}\par
\noindent\hspace*{6em}\texttt{"hexdump\_preview": bytes\_to\_hexdump(data[:64])\}}\par
\noindent\hspace*{3em}\texttt{report["files"].append(row)}\par
\noindent\hspace*{3em}\texttt{if claimed != actual:}\par
\noindent\hspace*{5em}\texttt{report["mismatches"].append(\{"path": path,}\par
\noindent\hspace*{7em}\texttt{"claimed\_type": claimed, "actual\_type": actual\})}\par
\noindent\hspace*{0em}\texttt{return report}\par
\end{tcolorbox}
\end{enumerate}
\end{workedbox}

\end{keepsection}
\begin{keepsection}\subsection{Task 10: Aggregate the forensic report and write the personalized justification}

\begin{workedbox}{Written justification \ensuremath{\cdot} use your own sample offsets \ensuremath{\cdot} worked solution}
\Needspace{14\baselineskip}
\begin{enumerate}
\item For the polyglot, name the file from your samples directory, list every format your validators accepted, and identify the byte evidence for each. Explain why both structures are valid rather than citing two coincidental marker strings.
\item For the appended payload, report the file type, the end-marker or structural rule used, the returned byte offset, total file length, and number of trailing bytes (file length minus the offset). Include a short hexdump around the boundary.
\item The guide intentionally does not publish a fixed filename or offset: this assignment seeds personalized examples. Run your own report and paste your actual results into your submission.
\item \Needspace{3\baselineskip}
\textbf{A starting point} for finding your two rows:
\begin{tcolorbox}[colback=Paperblue,colframe=Paperblue,boxrule=0pt,arc=0pt,left=7pt,right=7pt,top=5pt,bottom=5pt]\small\sloppy
\noindent\hspace*{0em}\texttt{\# Use your own generated files; never copy another learner's offsets.}\par
\noindent\hspace*{0em}\texttt{print(report["files"]) \# select the row with multiple valid types}\par
\noindent\hspace*{0em}\texttt{\# Then compare appended\_offsets[type] with that file's byte length.}\par
\end{tcolorbox}
\end{enumerate}
\end{workedbox}

\end{keepsection}

\vfill{\small Personal worked-solution guide based on the APDR Computer Exercise 1 task sheet, rewritten in my own words. Its ZIP-family requirement cannot be satisfied from the first 16 bytes alone; the page separates signature identification from archive inspection. The written justification must use each learner’s own seeded sample offsets. Not an official FAU publication. Made by \textbf{Thanate Amornjiraporn (Focus)}.}

\end{document}
